% Generated by tools/textbook/export_listings.js from tools/decks/mastery.js and tools/concepts/registry.json: do not edit.
\begin{chapstatus}{Stato di apprendimento di questo capitolo}
\masterychip{now}{DA COMPRENDERE OGGI}\enspace che cos'è l'informatica; l'architettura di von Neumann; RISC e un ISA aperto.\par\smallskip
\masterychip{use}{DA UTILIZZARE OGGI}\enspace il programma completo, gli strumenti di PRG32 (si usa, non si studia, in questo corso).\par\smallskip
\masterychip{later}{DA COMPRENDERE PIÙ AVANTI}\enspace l'algebra di Boole (Capitolo~\ref{ch:boole1}).
\end{chapstatus}
\begin{chapstatus}{Percorso del capitolo}
Alla fine del capitolo saprete:
\begin{enumerate}[nosep,leftmargin=2em]
\item definire l'informatica e spiegare perché non riguarda solo i computer
\item nominare le unità dell'architettura di von Neumann e la sua idea del programma memorizzato
\item dire che cosa distingue RISC da CISC e perché RISC-V è aperto
\end{enumerate}
Si appoggia su: solo l'installazione del Capitolo~\ref{ch:setup}. Prossimo passo: Capitolo~\ref{ch:information}.
\end{chapstatus}
